% --- 1. INFORMAZIONI GENERALI ---
\section{Informazioni Generali}
\begin{itemize}
	\item \textbf{Luogo/Piattaforma:} Server Discord\textsubscript{G}
	\item \textbf{Data:} 2026-05-05
	\item \textbf{Orario di inizio:} 14:30
	\item \textbf{Orario di fine:} 15:40
	\item \textbf{Responsabile\textsubscript{G}:} Anton Ricardo Rupi
	\item \textbf{Scriba:} Sofia De Blasi
\end{itemize}

\vspace{0.5cm}
\textbf{Partecipanti presenti:}
\begin{table}[ht]
	\centering
	\renewcommand{\arraystretch}{1.2}
	\begin{tabularx}{\textwidth}{X l}
		\toprule
		\textbf{Nome e Cognome} & \textbf{Matricola} \\
		\midrule
		Sofia De Blasi & 2111014 \\
		Roberto Mariano Doroftei & 2111031 \\
		Filippo Idiometri & 2035617 \\
		Francesco Marcon & 2110990 \\
		Armando Moda Scarati & 2082864 \\
		Anton Ricardo Rupi & 2054304 \\
		\bottomrule
	\end{tabularx}
\end{table}

\vspace{0.5cm}
\textbf{Partecipanti assenti:}
\begin{table}[ht]
	\centering
	\renewcommand{\arraystretch}{1.2}
	\begin{tabularx}{\textwidth}{X l}
		\toprule
		\textbf{Nome e Cognome} & \textbf{Matricola} \\
		\midrule
		Nessuno & - \\
		\bottomrule
	\end{tabularx}
\end{table}

% --- 2. ORDINE DEL GIORNO ---
\section{Ordine del Giorno}
Gli argomenti previsti per la riunione odierna sono:
\begin{enumerate}
	\item Sprint\textsubscript{G} review e retrospective,
	\item Pianificazione dello sprint successivo e assegnazione dei task\textsubscript{G} e ruoli.
\end{enumerate}

% --- 3. RESOCONTO E DISCUSSIONE ---
\section{Resoconto della Riunione}

\subsection{Sprint review\textsubscript{G} e retrospective}
Il gruppo si è confrontato sull'andamento dello sprint appena terminato. Questa settimana non ci sono stati particolari problemi. Riguardo il pdp, il file excel è stato diviso in più fogli organizzati per sprint, per ogni settimana il calcolo delle ore delle varie task è ora automatizzato. Il template del DdB è stato completato. Non sono emerse criticità per quanto riguarda l'AdR.

\subsection{Pianificazione dello sprint successivo e assegnazione dei task e ruoli}
Lo sprint è stato definito con durata dal 5 maggio al 12 maggio. I ruoli assegnati per questo sprint sono: 
Responsabile (Anton R. Rupi), Amministratore\textsubscript{G} (Sofia De Blasi), Verificatore\textsubscript{G} (Francesco Marcon), 
Analisti (Armando M. Scarati, Roberto M. Doroftei, Filippo Idiometri).
Inizia l'individuazione della struttura e contenuto del PdQ. Comincia anche la stesura del primo diario di bordo.
Come nei precedenti sprint, i ruoli di Programmatore\textsubscript{G} e Progettista\textsubscript{G} non sono ancora necessari e si è quindi 
deciso nuovamente di dedicare 3 membri al ruolo\textsubscript{G} dell'Analista\textsubscript{G}, in particolare Armando e Roberto riguardo all'AdR e Filippo riguardo l'individuazione delle tecnologie possibili da usare e valutare i pro e contro sia di quelle backend che frontend. Da questo sprint in poi la stesura del pdp verrà affidata all'amministratore.
% --- 4. DECISIONI PRESE ---
\section{Decisioni Prese}

\renewcommand{\arraystretch}{1.3}
\vspace{-0.3cm}
\noindent
\begin{longtable}{c p{12cm}}
	\toprule
	\textbf{Codice} & \textbf{Decisione} \\
	\midrule
	\endfirsthead

	\toprule
	\textbf{Codice} & \textbf{Decisione} \\
	\midrule
	\endhead

	\midrule
	\multicolumn{2}{r}{\textit{Continua nella pagina successiva}} \\
	\midrule
	\endfoot

	\bottomrule
	\endlastfoot

	\textbf{VI-12.1} & Iniziata l'individuazione delle tecnologie possibili e dei loro pro e contro \\[4pt]
	\textbf{VI-12.2} & Inizio la stilatura della struttura del PdQ\\[4pt]
	\textbf{VI-12.3} & La stesura del pdp è affidata all'amministratore da questo sprint in poi\\[4pt]
\end{longtable}

% --- 5. AZIONI (TODO) ---
\section{Azioni da Svolgere (To-Do)}

\renewcommand{\arraystretch}{1.3}
\begin{longtable}{c c p{4.5cm} p{3cm} l}
	\toprule
	\textbf{Codice} & \textbf{Rif. Decisione} & \textbf{Descrizione Task} & \textbf{Assegnatario} & \textbf{Scadenza} \\
	\midrule
	\endfirsthead

	\toprule
	\textbf{Codice} & \textbf{Rif. Decisione} & \textbf{Descrizione Task} & \textbf{Assegnatario} & \textbf{Scadenza} \\
	\midrule
	\endhead

	\midrule
	\multicolumn{5}{r}{\textit{Continua nella pagina successiva}} \\
	\midrule
	\endfoot

	\bottomrule
	\endlastfoot

	% Issue\textsubscript{G} #138 - Aggiornamento NdP
	\ghissue{138} & - & Aggiornare norme di progetto & Sofia De Blasi & 2026-05-12\\[4pt]
	
	% Issue #145 - doc. tecnico esterno analisi dei requisiti
	\ghissue{145} & VI-12.1 & Individuare tecnologie possibili e valutarne pro e contro & Filippo Idiometri & 2026-05-12 \\[4pt]
	
	% Issue #146 - Aggiornamento pdp
	\ghissue{146} & VI-12.3 & Stesura sprint 4 & Sofia De Blasi & 2026-05-12\\[4pt]
	
	% Issue #147 - scrivere diario di bordo
	\ghissue{147} & - & Redigere Diario di Bordo 1 (08/05) & Anton R. Rupi & 2026-05-08 \\[4pt]
	
	% Issue #148 - verbale interno 2026-05-05
	\ghissue{148} & - & Stesura verbale interno relativo a riunione tenutasi in data 2026-05-05 & Anton R. Rupi & 2026-05-12 \\[4pt]

	% Issue #149 - Aggiornamento AdR
	\ghissue{149} & - & Aggiornamento casi d'uso secondo le nuove convenzioni della macroarea 2 e 3 con estensioni per errori e approfondimento dell'Introduzione & Armando M. Scarati, Roberto M. Doroftei & 2026-05-12 \\[4pt]

	% Issue #151 - Iniziare la struttura del piano di qualifica
	\ghissue{151} & VI-12.2 & Individuare struttura e contenuto del Piano di Qualifica & Francesco Marcon & 2026-05-12 \\[4pt]

\end{longtable}

% --- 6. PROSSIMA RIUNIONE ---
\section{Prossima Riunione}
\begin{itemize}
	\item \textbf{Data:} 2026-05-12
	\item \textbf{Ora:} 14:30
	\item \textbf{OdG preliminare\textsubscript{G}:}
	\begin{itemize}
		\item Revisione\textsubscript{G} e verifica dei task del quinto sprint;
		\item Pianificazione del sesto sprint e assegnazione dei ruoli.
	\end{itemize}
\end{itemize}